%% kapa: exam=2 round=22 year=2011 period=2 subject=appraisal_theory subject_ko=감정평가이론 seq=1 number=1 points=40
%% source: Q-Net 감정평가사 기출문제 (cst00309) artlSeq=5040389 file=extracted/1008/5040389/2011년 제22회 감정평가사 2차시험 기출문제/2교시.hwp
%% conversion: kordoc 4.24.51 -> tools/gen2.py (auto); review=proofread-sonnet; flags=none
\begin{longquestion}{【문제 1】}
부동산의 가치는 여러 가지 요인에 의해 영향을 받기 때문에 감정평가사는 대상부동산의 개별적인 특성뿐만 아니라 정부의 정책과 부동산시장변화에 대해서도 이해할 필요가 있는 바, 다음의 물음에 답하시오. (40점) \par
\medskip\noindent (1) 최근 전력난을 완화하기 위한 초고압 송전선로 설치가 빈번하게 발생하고 있으며 이를 둘러싼 이해관계자들의 갈등도 증폭되고 있는데, 이와 관련된 선하지 보상평가방법과 송전선로 설치에 따른 ‘보상되지 않는 손실’에 대해 설명하시오. (15점) \par
\medskip\noindent (2) 최근 수익형 부동산에 대한 관심이 확산되고 있는데 수익형 부동산의 특징과 그 가격형성원리에 대해 설명하시오. (15점) \par
\medskip\noindent (3) 수익형 부동산의 평가방법에 대해 설명하시오. (10점) \par
\end{longquestion}
